\scanpage{192}
\mspage{2.}
définit un morphisme de sites en sens inverse).

Axiomes d'une topologie : a) Si une famille est \unsure{majorisée|majorée} par une famille couvrante, elle est couvrante : trivial

b) Un \unsure{isom|somme} est couvrant : trivial \quad c) stabilité par chgt. de base : on est ramené : des $X_i-X$ qui sont des monomorphismes, grâce : b).
\diagram{\begin{tikzcd}[ampersand replacement=\&]
X_i\arrow[d] \& X'_i\arrow[l]\arrow[d]\\
X \& X'\arrow[l]
\end{tikzcd}}
mais alors (mettre en lemme) $r$ commute aux produits fibrés $X_i\times_X X'$, et on gagne

d) Notion locale : supposons qu'on ait des $X'_\alpha-X$ couvrant $X$, tels que $\forall\alpha$, les $X_i\times_\alpha X'_\alpha=X'_{i\alpha}$ couvrent $X_\alpha$. Alors les $X_i$ couvrent $X$. Ici encore, on est ramené au cas où les $X_i\to X$ sont des monomorphismes, et on gagne.
\diagram{\begin{tikzcd}[ampersand replacement=\&]
X_i\arrow[d,two heads] \& X'_{i\alpha}\arrow[d,two heads]\\
Y_i\arrow[d,hook] \& Y'_{i\alpha}\arrow[d,hook]\\
X \& X'_\alpha\arrow[l]
\end{tikzcd}}
(Question : y a-t-il une référence à SGA 4 : image inverse d'une topologie ? Mais attention au fait que $r$ n'est pas exact à g.!). On n'a utilisé pour l'instant que R1. R3 \deleted{implique}\inserted{signifie} que la topologie est \deleted{\unsure{plus}}\unsure{moins} fine que la canonique, R4 que \inserted{les objets de $\mathscr R_{S'}$ sont loc. cts finis}. L'assertion sur les changements de base est triviale.

\uline{Corollaire} Prenant les topos \deleted{\unsure{assoc|associés aux}} $\widetilde{\mathscr R}_{S'}$ associés aux $\mathscr R_{S'}$, on trouve sur $Et(S)$ un \uline{champ en topos}. $\mathscr R_{S'}$ ``sont'' la sous-catégorie des objets loc. constants finis de $\widetilde{\mathscr R}_{S'}$

\deleted{Soit $X\to Y$ un morphisme de $\mathscr R_{S'}$. Alors $Y=\coprod_{n\in\mathbb N}Y_n$, \inserted{(somme \uuline{effective}} avec $X\times_Y Y_n\to Y_n$ \uline{de rang} $n$. \deleted{\unsure{Quelle}} \inserted{La famille} \deleted{\unsure{somme}} $Y_n\to Y$ est elle \uuline{couvrante} ?? \inserted{En d'autres termes,} Les $r(Y_n)\to r(Y)$ couvrent-ils ??}
% archive sheet number: 191
